\subsection{Parallel algorithms}

\begin{frame}[t,fragile]{Algorithms and execution policy}
\begin{itemize}
  \item Many algorithms accept an \textmark{execution policy}:
\begin{lstlisting}
std::vector<double> v = get_values();
auto n1 = std::count_if(v.begin(), v.end(), [](double x) { return x>0; });
auto n2 = std::count_if(std::execution::seq, v.begin(), v.end(), [](double x) { return x>0; });
auto n3 = std::count_if(std::execution::par, v.begin(), v.end(), [](double x) { return x>0; });
\end{lstlisting}

  \mode<presentation>{\vfill\pause}
  \item Counting values that satisfy a predicate with \cppcode{count_if()}:
\begin{lstlisting}
#include <algorithm>
#include <execution>
#include <print>
#include <vector>
#include "image.hpp"

int main() {
  std::vector<image> v = load_images("img.dat");
  auto n = std::count_if(std::execution::par, v.begin(), v.end(), 
      [](image const & im) { return not im.is_valid(); });
  std::println("Invalid images = {}", n);
}
\end{lstlisting}

\end{itemize}
\end{frame}

\begin{frame}[t,fragile]{Execution policies}
\begin{itemize}
  \item \cppcode{std::execution::seq}:
    \begin{itemize}
      \item Execution is sequenced as if on a single thread (no parallelism or vectorization assumed).
    \end{itemize}
  \item \cppcode{std::execution::par}:
    \begin{itemize}
      \item May execute iterations in parallel, but each iteration is sequential (no vectorization assumed).
    \end{itemize}
  \item \cppcode{std::execution::par_unseq}:
    \begin{itemize}
      \item May execute iterations in parallel and unsequenced; operations may be vectorized.
    \end{itemize}
  \item \cppcode{std::execution::unseq}:
    \begin{itemize}
      \item Executes unsequenced within a single thread context; operations may be vectorized.
    \end{itemize}
\end{itemize}
\end{frame}

\begin{frame}[t,fragile]{Parallel algorithms and exceptions}
\begin{itemize}
  \item If an exception is thrown by element access, predicate, or operation during an algorithm invoked with a \cppcode{std::execution} policy, \cppcode{std::terminate()} is called.

  \mode<presentation>{\vfill\pause}
  \item This includes \cppcode{std::execution::seq}, even though it is not parallel.
    \begin{itemize}
      \item Does not apply to non-\cppcode{std::execution}-based algorithms.
    \end{itemize}

  \mode<presentation>{\vfill\pause}
  \item An algorithm may still throw \cppcode{std::bad_alloc} if memory allocation fails, even when using a parallel execution policy.
  
 
  \mode<presentation>{\vfill\pause}
  \item \textmark{Rationale}: Parallel algorithms may execute work on multiple threads or in vectorized/unsequenced contexts, and the execution-policy overloads do not support propagating exceptions from user code to the caller.
    \begin{itemize}
      \item Calling \cppcode{std::terminate()} keeps the semantics simple for implementations and preserves optimization freedom for parallel and vectorized execution.
    \end{itemize}
\end{itemize}
\end{frame}